\documentclass[11pt]{article}
\usepackage{varnernotes}

\newcommand{\R}{\mathbb{R}}
\newcommand{\D}{\mathcal{D}}

\begin{document}

\lecturenotes{1b}{Time Value of Money, Abstract Assets, and Net Present Value}{Fall 2026}{August 27, 2026}{Jeffrey D. Varner}

\learningobjectives
  {Convert dated dollars}
  {Explain why cash flows at different dates require common units and derive discrete and continuous discount factors.}
  {Asset valuation}
  {Represent a project, product, or security as an abstract asset and compute its net present value.}
  {Interpret the benchmark}
  {Explain what the sign of NPV does---and does not---say under a declared discount curve and cash-flow model.}

\section{Why Does Time Have a Price?}

The time value of money begins with a simple asymmetry: one dollar in hand today and one dollar promised next year carry the same currency label, but they do not offer the same set of choices. Today's dollar can be consumed, held as a reserve, or invested. The future dollar cannot be used until it arrives, and its purchasing power and delivery are uncertain.

This idea predates modern financial markets. Irving Fisher organized it around two forces: impatience to spend income and the opportunity to invest it. Modern valuation adds inflation and risk to the same basic story. A discount rate is therefore not merely a computational input; it records the return forgone by waiting and, when appropriate, compensation for uncertainty.

L1a distinguished policy rates, bank funding rates, and asset-specific valuation rates. Here the discount rate is a declared model input: it describes the conversion between dated dollars under a stated benchmark and compounding convention. Monetary policy may influence that benchmark, but it does not uniquely determine the correct valuation rate for every asset.

\section{Discounting as a Change of Units}

A useful currency model treats money at different dates as different currencies: $\mathrm{USD}_0,\mathrm{USD}_1,\ldots$. Just as euros and U.S. dollars cannot be added before conversion, dated dollars cannot be added until they share a valuation date. The quoted yield specifies the conversion across time, and the accumulation and discount factors perform it. An abstract asset then provides a flexible model of the economic value of a process, good, or idea as dated signed cash flows (\figref{abstract-asset}).

\begin{figure}[ht]
  \centering
  % Shared with the companion notebooks -- see figs/abstract-asset-timeline.tex.
  % Do not inline a copy here; edit the shared file so both artifacts stay in step.
  % Drawing styles come from templates/vnflow.sty via varnernotes.sty.
  % Abstract-asset timeline -- SHARED SOURCE.
%
% \input by the standalone wrapper below; do not duplicate the picture.
% Rebuild the PDF/SVG for the notebook and slides with `make` in this directory.
%
% Drawing grammar and the CHEME flow convention come from templates/vnflow.sty:
%   arrow LEAVING a time-point unit  = cash paid
%   arrow ENTERING a time-point unit = cash received
% The signed net cash flow \bar c_j already carries its sign, so labels are
% written WITHOUT an extra minus; the arrow direction conveys direction.
%
% Indexing matches the notebooks and L2a: j = 0,1,...,N with N = nT.
\begin{tikzpicture}[x=2.25cm,y=1cm,>=Latex,font=\sffamily\small]

  \draw[vn brace] (0.35,1.65) -- (6.15,1.65)
    node[midway,above=7pt,text=vnink]
    {\bfseries abstract asset $\mathcal{A}=\{\bar c_0,\bar c_1,\ldots,\bar c_N\}$};

  \draw[vn timeline] (0,0) -- (6.65,0) node[right] {time};
  \foreach \x/\lab/\unit in {0.35/0/u0,1.75/1/u1,3.15/2/u2,6.0/N/uN}{
    \node[vn unit] (\unit) at (\x,0) {};
    \node[below=4pt] at (\x,0) {$j=\lab$};
  }
  \node at (4.55,0.02) {$\cdots$};

  % j = 0: arrow leaves the unit -> cash paid (the purchase)
  \draw[vn paid] (u0.north) -- (0.35,1.05);
  \node[right,text=vncarnelian] at (0.35,0.62) {$\bar c_0$};
  \node[above,vn annotation] at (0.35,1.05) {cash paid};

  % j = 1: arrow enters the unit -> cash received
  \draw[vn received] (1.75,1.05) -- (u1.north);
  \node[right,text=vnlink] at (1.75,0.62) {$\bar c_1$};

  % j = 2: arrow leaves the unit -> cash paid
  \draw[vn paid] (u2.north) -- (3.15,1.05);
  \node[right,text=vncarnelian] at (3.15,0.62) {$\bar c_2$};

  % j = N: arrow enters the unit -> cash received (the sale)
  \draw[vn received] (6.0,1.05) -- (uN.north);
  \node[right,text=vnlink] at (6.0,0.62) {$\bar c_N$};
  \node[above,vn annotation] at (6.0,1.05) {cash received};

  % Transport every dated flow to the common valuation date
  \draw[vn transport] (6.0,-0.94) -- (0.50,-0.94)
    node[midway,fill=white,inner sep=3pt,text=vnlink]
    {$\bar c_j\mapsto\mathcal{D}_{j,0}^{-1}(y;n)\,\bar c_j$};
  \node[vn callout,anchor=north west] at (0.35,-1.27)
    {common units: time-0 dollars};
  \node[vn annotation,anchor=west] at (3.35,-1.57)
    {discount future cash flows to the valuation date};
\end{tikzpicture}

  \caption{The abstract asset $\mathcal{A}$ has a signed net cash flow $\bar c_j$ at each index $j\in\{0,1,\ldots,N\}$, where $N=nT$ is the final cash-flow index. An arrow entering a timeline marker is money received; an arrow leaving it is money paid. The sign is already included in $\bar c_j$, so the labels need no extra minus sign. Discount each future cash flow back to the common valuation date $j=0$.}
  \label{fig:abstract-asset}
\end{figure}

Start with one year of annual compounding. Suppose the initial principal $P$ dollars earns the nominal annual yield $y$ for one year. The resulting year-end future value $F$ is given by \eqnref{annual-compounding}:
\begin{equation}
  \label{eq:annual-compounding}
  F=(1+y)P.
\end{equation}
With annual compounding, the quoted yield $y$ produces the forward conversion factor $1+y$ and the backward factor $(1+y)^{-1}$. These two directions define a reversible conversion between dated currencies (\defnref{discount-factor}) and let us express every cash flow in \figref{abstract-asset} at a common date.

\begin{definition}{Accumulation and discount factors}{discount-factor}
Let time $0$ be the valuation date, let the compounding-period index $j\geq 0$ identify a future date, and let $\mathrm{USD}_0$ and $\mathrm{USD}_j$ denote dollars measured at those respective dates. For a quoted nominal annual yield $y$ and a stated compounding convention, the \emph{forward accumulation factor} $\D_{j,0}(y)$ is the multiplicative conversion factor from $\mathrm{USD}_0$ to $\mathrm{USD}_j$. If the present amount $P_0$ is measured in $\mathrm{USD}_0$ and the equivalent future amount $F_j$ is measured in $\mathrm{USD}_j$, then $F_j=\D_{j,0}(y)P_0$ and $P_0=\D_{j,0}^{-1}(y)F_j$, where $\D_{j,0}^{-1}(y):=1/\D_{j,0}(y)$ is the \emph{discount factor}. A valid pair satisfies $\D_{0,0}=1$.

Throughout this course, $\D$ without an inverse always accumulates forward and $\D^{-1}$ always discounts back. Many finance texts use the phrase \emph{discount factor} for both objects; the formulas are equivalent once the convention is stated, but we keep the two names distinct.
\end{definition}

The periodic-compounding formula is a model rather than an algebraic identity. It assumes that the nominal annual yield $y$ remains constant; the positive integer compounding frequency $n$ divides each year into $n$ equal intervals; interest is credited at the end of every interval at the proportional rate $y/n$; credited interest remains invested; and the investment horizon $T$ years contains the whole number $nT$ of intervals with no intervening deposits or withdrawals. That last assumption governs the accumulation of a single principal; it places no restriction on the cash-flow streams we discount in \S3, which may have flows at any date.

Let the account value $V_k$ denote the value immediately after interval $k$, and let the initial principal be $V_0=P$. Applying the same one-interval update repeatedly gives the periodic accumulation factor:
\begin{align*}
  V_0
    &= P,
    && \text{\color{vnmuted}\footnotesize initial deposit} \\
  V_1
    &= \left(1+\frac{y}{n}\right)V_0
     = P\left(1+\frac{y}{n}\right),
    && \text{\color{vnmuted}\footnotesize apply the interval multiplier once} \\
  V_2
    &= \left(1+\frac{y}{n}\right)V_1
     = P\left(1+\frac{y}{n}\right)^2,
    && \text{\color{vnmuted}\footnotesize apply it a second time} \\
  &\ \vdots
    && \text{\color{vnmuted}\footnotesize repeat the same update} \\
  V_{j}
    &= P\left(1+\frac{y}{n}\right)^{j},
    && \text{\color{vnmuted}\footnotesize after $j$ intervals} \\
  \D_{j,0}(y;n)
    &:= \frac{V_{j}}{V_0}
     = \left(1+\frac{y}{n}\right)^{j},
    && \text{\color{vnmuted}\footnotesize normalize by the initial value.}
\end{align*}
The exponent $j$ counts the compounding intervals and the base $1+y/n$ is the one-interval multiplier, valid whenever $1+y/n>0$. A horizon of $T$ years that lands on the compounding grid corresponds to $j=nT$, so we may write the same factor against elapsed time as $\D_{T,0}(y;n)=(1+y/n)^{nT}$; the first subscript then carries years rather than a period count. As the compounding frequency $n$ increases, this factor approaches the continuous-compounding limit (\thmref{continuous-limit}).

\begin{theorem}{Continuous-compounding limit}{continuous-limit}
Let the nominal annual yield be the constant $y\in\R$, where $\R$ denotes the set of real numbers; let the investment horizon be $T\geq 0$ years; and let the positive integer $n$ denote the number of compounding intervals per year, with $1+y/n>0$. As the compounding frequency $n$ increases without bound, the discrete accumulation factor converges to its continuous counterpart, where $e$ denotes Euler's number, as stated by \eqnref{continuous-limit-equation}:
\begin{equation}
  \label{eq:continuous-limit-equation}
  \lim_{n\to\infty}\left(1+\frac{y}{n}\right)^{nT}=e^{yT}.
\end{equation}
Consequently, a continuously compounded rate $g$ has forward accumulation factor $\D_{T,0}(g)=e^{gT}$ and discount factor $\D_{T,0}^{-1}(g)=e^{-gT}$.
\end{theorem}

\begin{derivation}
Apply the natural logarithm $\log$ to the discrete accumulation factor $\D_{T,0}(y;n)$ and use the Taylor expansion $\log(1+x)=x-x^2/2+O(x^3)$, where $x=y/n$ and $O(x^3)$ denotes a remainder bounded in magnitude by $C\lvert x\rvert^3$ near $x=0$ for some constant $C>0$:
\begin{align*}
  \log \D_{T,0}(y;n)
  &=nT\log\left(1+\frac{y}{n}\right) \\
  &=yT-\frac{y^2T}{2n}+O(n^{-2}).
\end{align*}
Here $O(n^{-2})$ denotes terms bounded in magnitude by $C'/n^2$ for sufficiently large compounding frequency $n$ and some constant $C'>0$. These correction terms vanish as $n\to\infty$; exponentiating gives the result. The expansion also shows how quickly discrete compounding approaches the continuous convention.
\end{derivation}

\thmref{continuous-limit} is a statement about a limit, and it is easy to over-read. It does \emph{not} say that a quoted nominal yield is already continuously compounded. For a \emph{fixed} finite $n$, the exact continuously compounded equivalent of $y$ is given by \eqnref{continuous-equivalent}:
\begin{equation}
  \label{eq:continuous-equivalent}
  g_y:=n\log\left(1+\frac{y}{n}\right),
  \qquad\text{so that}\qquad
  \left(1+\frac{y}{n}\right)^{nT}=e^{g_yT}.
\end{equation}
Equation~\eqref{eq:continuous-equivalent} is an identity, not an approximation: $y$ with $n$ compounding intervals and $g_y$ compounded continuously are two spellings of the same accumulation rule. The limit in \eqnref{continuous-limit-equation} instead compares $y$ against a \emph{different} rate as $n$ grows, and the gap it closes is $O(y^2T/2n)$. Over a long horizon at a large yield that gap is not small: the companion notebook tabulates it and shows the leading-order estimate failing badly at $n=1$.

Writing $g$ for a continuously compounded rate also fixes a unit convention we use for the rest of the course: $g$ carries units of inverse time, so the dimensionless log return accumulated over an interval $\Delta t$ is $r=(\Delta t)\,g$. L2a uses this to compare a bill's cumulative gain $r_B$ with its annualized growth $g_B$, and L3a uses it for equity log returns.

The companion notebook turns these expressions into computed factors. Use it to compare discrete and continuous compounding numerically and to see how quickly the two conventions converge.

\notebooklink
  {L1b: Time Value of Money}
  {https://github.com/varnerlab/CHEME-5660-CourseRepository-Fall-2026/blob/main/lectures/week-1/L1b/CHEME-5660-L1b-Lecture-TimeValueMoney-Fall-2026.ipynb}
  {Build discrete and continuous discount factors and compare the two compounding conventions numerically.}

\section{Abstract Assets and Net Present Value}

With the conversion rule in hand, we can return to the abstract asset. Let $\bar c_j$ be the signed net cash flow in $\mathrm{USD}_j$: receipts are positive and payments are negative. An asset with $N+1$ cash-flow dates is represented by $\{\bar c_j\}_{j=0}^{N}$. NPV follows a simple discipline---\emph{convert, then add}: multiply each $\bar c_j$ by $\D_{j,0}^{-1}(y;n)$ to express it in $\mathrm{USD}_0$, then sum the converted values (\defnref{net-present-value}).

\begin{definition}{Net present value}{net-present-value}
Let time $0$ be the valuation date, let the nonnegative integer $N$ be the terminal cash-flow index, and let each integer $j\in\{0,1,\ldots,N\}$ identify a cash-flow date. Let the signed net cash flow $\bar c_j$ be measured in $\mathrm{USD}_j$, and let the accumulation factor $\D_{j,0}(y;n)$ convert $\mathrm{USD}_0$ to $\mathrm{USD}_j$ under the nominal annual yield $y$ compounded $n$ times per year. Its reciprocal $\D_{j,0}^{-1}(y;n):=1/\D_{j,0}(y;n)$ converts the cash flow back to time-0 dollars. The net present value at time $0$ is given by \eqnref{npv}:
\begin{equation}
  \label{eq:npv}
  \operatorname{NPV}_0(y)=\sum_{j=0}^{N}\D_{j,0}^{-1}(y;n)\,\bar c_j
  =\left\langle\D_{\star,0}^{-1}(y;n),\bar{\mathbf c}\right\rangle,
\end{equation}
where the $\star$ subscript denotes the ordered $(N+1)$-vector over $j=0,\ldots,N$. Every term is measured in time-0 dollars, so the sum produces one NPV value.
\end{definition}

The sign of NPV compares the asset with the benchmark represented by the discount curve. The result depends on the stated cash flows and yield. It does not guarantee a realized profit or better performance than another investment:

\textbf{Positive NPV ($\operatorname{NPV}>0$).}\enspace Discounted receipts exceed discounted payments. Under the stated cash-flow model and yield, the asset is worth more than the benchmark it is being priced against.

\textbf{Zero NPV ($\operatorname{NPV}=0$).}\enspace Discounted receipts equal discounted payments, so the asset is fairly priced relative to that benchmark. Zero NPV does not mean zero nominal gain and does not mean zero risk.

\textbf{Negative NPV ($\operatorname{NPV}<0$).}\enspace Discounted payments exceed discounted receipts, so the asset does not meet the required return used in the calculation. This does not by itself imply a nominal cash loss: undiscounted receipts may still exceed undiscounted payments once time is ignored.

The worked-example notebook applies the abstract-asset model to a vehicle purchase. It constructs the dated cash flows, builds the accumulation and discount factors, and computes the time-0 NPV as a scalar product.

\notebooklink
  {L1b: Net Present Value Worked Example}
  {https://github.com/varnerlab/CHEME-5660-CourseRepository-Fall-2026/blob/main/lectures/week-1/L1b/CHEME-5660-L1b-NetPresentValue-Fall-2026-WorkedExample.ipynb}
  {Represent a vehicle purchase as an abstract asset, build its discount factors, and compute the time-0 NPV.}

\thispagestyle{fancy}
\keytakeaways
  {In this lecture, we placed dated cash flows in common time-0 units, derived discount factors, and defined NPV for an abstract asset.}
  {Dated dollars require conversion}
  {Cash flows at different dates cannot be added economically until a discount factor expresses them at a common valuation date.}
  {Compounding is a convention}
  {A nominal yield $y$ compounded $n$ times per year and its continuous equivalent $g_y=n\log(1+y/n)$ are the same rule written two ways; neither should be used without stating its assumptions.}
  {Structure before identity}
  {An abstract asset describes any investment as a dated sequence of payments and receipts.}
  {Next, use the companion notebooks to build these factors numerically and compute an NPV end to end.}

\begin{readinglist}
  \item Irving Fisher, \href{https://oll.libertyfund.org/titles/fisher-the-theory-of-interest}{\emph{The Theory of Interest}} (1930), Chapters I--III, for the classical account of impatience and investment opportunity.
  \item Jonathan Berk and Peter DeMarzo, \href{https://www.pearson.com/en-gb/subject-catalog/p/Berk-Corporate-Finance-Global-Edition-5th-Edition/P200000003710}{\emph{Corporate Finance}}, 5th ed. (2020), Chapters 3--4, for present value and investment decision rules.
  \item Frank J. Fabozzi and Francesco A. Fabozzi, \href{https://mitpress.mit.edu/9780262046275/bond-markets-analysis-and-strategies/}{\emph{Bond Markets, Analysis, and Strategies}}, 10th ed. (2021), for discount factors and fixed-income conventions.
\end{readinglist}

\vfill
{\sffamily\footnotesize\color{vnmuted}\textbf{Educational use and risk notice.} These notes are for instruction only and do not constitute investment advice. Financial decisions involve risk, and model outputs depend on assumptions.}

\end{document}
